\chapter{ADCS – Control y Determinación de Actitud}

\section{Introducción al Control de Posición Orbital}
El módulo de control de actitud (ADCS, \textit{Attitude Determination and Control Subsystem}) tiene como fin estabilizar giros remanentes (detumbling) arrastrados tras la eyección y fijar un referencial estable respecto a vectores terrestres. Operando de forma acoplada con el OBC, resulta crítico para la fase nominal de la misión (apuntamiento fotográfico) y del Modo de Supervivencia.

\section{Esquema Pasivo y Dinámico de Actitud}

\subsection{Estabilización Magnética Pasiva y Materiales}
El primer nivel de respuesta ante control de rotación de bajo costo es un diseño pasivo térmicamente tolerante:
\begin{itemize}
    \item \textbf{Cuerpo Magnético Guía:} Uno o más magnetos permanentes (ej. Neodimio N45) alojados directamente en el chasis o la tarjeta dedicada ADCS, que obligan al satélite a acoplarse y arrastrarse por las líneas estáticas magnéticas formadas por los polos Norte/Sur del planeta (el eje de satélite busca alinearse latitudinalmente).
    \item \textbf{Barras disipativas (Hysteresis):} Cintas mecánicas perpendiculares a la fuerza (como \textit{Permanorm}) que inducen la fricción magnética generando amortiguamiento al esparcir la energía rotacional descontrolada (\textit{detumbling}) del CubeSat convirtiéndola pasivamente en calor latente.
\end{itemize}

\subsection{Gestión Activa y Módulos de Sensores}
Para lograr la precisión necesaria para apuntar los paneles solares continuamente (Sun-Pointing) y asistir microscopía y ruteos de telecomunicación:
\begin{itemize}
    \item \textbf{Determinación:} Múltiples magnetómetros en los ejes cartesianos interactúan con Fotodiodos LDR sensando flujos solares orbitales. La matriz es alimentada al bus general de medición del ADCS leyendo con precisión angular de alta velocidad el punto orbital actual.
    \item \textbf{Accionadores (Actuadores Dinámicos):} Un conjunto de tres a más \textit{Magnetorquers} (bobinas embebidas en las PCB con polarización reversible o tubos cilíndricos de cobre) provocan momentos dipolares artificiales contra el campo terrestre modificando voluntariamente la latitud para encauzar a la luz en los puertos fotovoltaicos, reaccionando estrictamente en el estado \textit{Modo de Seguridad} a órdenes directas inyectadas por la OBC.
\end{itemize}
